% ═══════════════════════════════════════════════════════════════════
% §7  지열원 히트펌프 급탕기 (동적 모델)
% Source: GroundSourceHeatPumpBoiler.py
% ═══════════════════════════════════════════════════════════════════
\section{지열원 히트펌프 급탕기}
\label{sec:gshpb}

본 절은 \texttt{GroundSourceHeatPumpBoiler.py}의 동적 GSHPB 모델을 기술한다. ASHPB(\ref{sec:ashpb}절)의 증기압축 프레임워크를 공유하되, 공기열원 증발기를 지중 결합 보어홀 열교환기로 대체한다.

\subsection{가정}
\begin{enumerate}
  \item 지중 온도는 깊이 방향으로 균일하고 일정하다 (비교란 지중 온도 $T_g$).
  \item 보어홀 열교환기는 유한선원(FLS) g-함수로 모델링한다.
  \item 지중 루프 유체 물성은 평균 브라인 온도에서 평가한다.
  \item 시스템은 CoolProp 기반 물리적 물성치를 통해 명시적 증기압축 사이클로 해석되며, 응축기 모델은 열역학적 밸런스를 고려한 침수 코일 방식으로 연산된다.
\end{enumerate}

\subsection{지중 결합 증발기}

지중으로부터의 채열량:
\begin{equation}\label{eq:gshp-Qground}
  Q_\text{ground} = \frac{T_g - T_\text{brine,in}}{R_b + R_g(t)}
\end{equation}
여기서 $R_b$는 보어홀 열저항, $R_g(t)$는 g-함수(식~\ref{eq:fls-g})로부터 산출되는 시간 의존 지중 열저항이다:
\begin{equation}\label{eq:gshp-Rg}
  R_g(t) = \frac{g(t)}{2\pi k_s H}
\end{equation}
$k_s$는 지중 열전도도, $H$는 보어홀 깊이이다.

\subsection{명시적 증기압축 역학 모델 (Explicit Vapor-Compression Cycle)}

단순 Carnot 기반의 경험적 상관식 적용을 탈피하여, 본 모듈은 CoolProp 라이브러리와 연계한 명시적 증기압축 사이클 모델을 해석한다. 증발기 측의 열전달 성능과 냉매의 증발 온도 하강분($\Delta T_{\text{evap}}$)을 동적으로 결정하기 위해 Brent-최적화 기법을 활용한다. 매 시간 단계에서 열역학적 엔탈피 밸런스를 풀이하여 냉매 질량 유량, 압축기 일량, 응축열량을 정밀 도출하며, 소모 전력 $E_{\text{comp}}$와 성능 계수 $\text{COP}_{\text{sys}}$는 다음과 같이 산출된다:
\begin{equation}\label{eq:gshp-w-comp}
  E_{\text{comp}} = \dot{m}_{\text{ref}} (h_{\text{comp,out}} - h_{\text{comp,in}})
\end{equation}
\begin{equation}\label{eq:gshp-COP}
  \text{COP}_\text{sys} = \frac{Q_\text{cond}}{E_{\text{comp}} + E_{\text{pump,brine}}}
\end{equation}
여기서 $\dot{m}_{\text{ref}}$는 냉매 질량 유량이고, $h_{\text{comp,out}}$는 내부 등엔트로피 효율($\eta_{\text{isen}}$)을 반영하여 결정된다. 이 역학적 해는 사용자 설정 하중과 동적 지중열원 온도를 양립하는 유일한 열역학적 작동점으로 수렴한다.

\subsection{브라인 루프 에너지 수지}

\begin{equation}\label{eq:gshp-brine}
  T_\text{brine,out} = T_\text{brine,in} - \frac{Q_\text{evap}}{\dot{m}_\text{brine}\,c_\text{brine}}
\end{equation}

\subsection{엑서지 분석}

엑서지 수지 구조는 ASHPB(\ref{sec:ashpb}절)와 유사하며, 실외기가 지중 결합 루프로 대체된다:
\begin{equation}\label{eq:gshp-Xground}
  X_\text{ground} = Q_\text{ground}\left(1 - \frac{T_0}{T_g}\right)
\end{equation}

전체 엑서지 효율:
\begin{equation}\label{eq:gshp-etaX}
  \eta_{X,\text{sys}} = \frac{X_\text{w,out} - X_\text{w,in}}{X_\text{cmp} + X_\text{pump,brine}}
\end{equation}
